% !TeX root = ../../Book1.tex
\section{波兰空间与相容的完备度量}
\label{b1:sec:A105}

第\ref{b1:sec:1204}节的波兰空间结合两种优势：可数基使测试条件可以可数化，某个相容度量的完备性使 Cauchy 逼近留在空间内。这里补充这一类空间在构造下的稳定性。由于定义中的完备度量只要求存在，我们可以改变距离的数值，只要不改变开集。

\subsection{开子集为什么仍可完备度量化}
\label{b1:subsec:A10501}

\begin{proposition}\label{b1:prop:app-polish-operations}
波兰空间的闭子空间、可数乘积仍是波兰空间。若 \(X\) 为波兰空间且 \(G\subset X\) 是可数个开集之交，则 \(G\) 也是波兰空间；特别地，开子空间也是。
\end{proposition}
\begin{proof}
闭子空间的断言由完备性继承与可分度量子空间的可分性得到。可数乘积用\cref{b1:thm:app-countable-product}，对每个因子先选相容完备度量。

令 \(G=\bigcap_{j\ge1}U_j\)，在 \(X\) 上固定相容完备度量 \(d\)。若 \(G\) 为空，断言直接成立。对非空补集 \(X\setminus U_j\)，在 \(G\) 上设
\[
 h_j(x)=\frac1{d(x,X\setminus U_j)};
\]
若补集为空，令 \(h_j=0\)。每个分母在 \(G\) 上为正，故 \(h_j\) 连续。考虑图映射
\[
 \Phi:G\longrightarrow X\times\mathbb R^{\mathbb N},
 \quad x\longmapsto\bigl(x,h_1(x),h_2(x),\ldots\bigr).
\]
它与第一坐标投影互逆，因此是到其像的同胚。下面证明像闭。

设图上序列 \(\Phi(x_n)\) 在乘积中趋向 \((x,t_1,t_2,\ldots)\)。若 \(x\notin U_j\) 且该补集非空，则 \(d(x_n,X\setminus U_j)\le d(x_n,x)\rightarrow0\)，从而 \(h_j(x_n)\rightarrow+\infty\)，与 \(t_j\) 有限矛盾。所以 \(x\in G\)，各坐标连续性再给出 \(t_j=h_j(x)\)。乘积可度量，序列判据足以说明像闭。该乘积有相容完备度量，闭子空间因而完备；把它拉回 \(G\)，便得与原子空间拓扑相容的完备度量。可分性已由 \(G\subset X\) 继承。
\end{proof}

这个构造并没有把边界点补进 \(G\)。恰恰相反，倒数距离使趋向被删边界的序列在某个新坐标上逃向无穷，从而不再是新度量下的 Cauchy 列。可把拉回度量具体写成
\[
 \rho(x,y)=\min\{1,d(x,y)\}
 +\sum_{j=1}^{\infty}2^{-j}\min\{1,|h_j(x)-h_j(y)|\}.
\]
图闭的论证说明它完备，而有限前段加小尾和的论证说明它不改变拓扑。这一形式与 Fremlin《Measure Theory》卷四\cite[4A2M、4A2Q]{Fremlin2016Measure}中的完备度量构造相联系。

\subsection{例子与不能省去的条件}
\label{b1:subsec:A10502}

欧氏空间的开集、闭集都是波兰空间。无理数集也是，因为它是可数个开集 \(\mathbb R\setminus\{q\}\)、\(q\in\mathbb Q\)，的交。无理数集的通常距离并不完备，例如无理数列可以趋向有理数；上一构造给出的是另一相容完备度量。

有理数集却不是波兰空间。若其通常拓扑可由某个完备度量给出，则\cref{b1:thm:baire}适用。每个有理数单点闭且没有内点，所以 \(\mathbb Q\) 是可数个无处稠密闭集的并，与非空完备度量空间的 Baire 性矛盾。这里无处稠密性只依赖拓扑，因此换成任何相容度量也不能消除矛盾。

由此还得到一个值得保留的区别：波兰空间的任意子集不一定波兰，即使它是 Borel 集。\(\mathbb Q\) 是 \(\mathbb R\) 中的可数并闭集，却不波兰。相反，可数交开集的结论有上面的明确构造作为保证，不能把“Borel”三个字直接替换进去。

\begin{proposition}\label{b1:prop:app-polish-functions}
若 \(K\) 是紧度量空间，则实连续函数空间 \(C(K)\) 在一致范数下为可分 Banach 空间，因而是波兰空间。
\end{proposition}
\begin{proof}
一致 Cauchy 列逐点有极限，并由共同的 Cauchy 估计一致收敛；一致极限连续，所以 \(C(K)\) 完备。以下证可分，空 \(K\) 的情形直接成立。

固定可数稠密集 \(D\subset K\)。考虑所有有限个 \(a_i\in D\)、正有理数 \(r\)、有理数 \(q_i\) 所给出的函数
\[
 g(x)=\frac{\sum_{i=1}^N q_i\max\{0,r-d(x,a_i)\}}
 {\sum_{i=1}^N\max\{0,r-d(x,a_i)\}},
\]
只保留分母在整个 \(K\) 上为正的选择。它们组成可数族。给定 \(f\in C(K)\) 与 \(\varepsilon>0\)，一致连续性允许选正有理数 \(r\)，使 \(d(x,y)<r\) 时 \(|f(x)-f(y)|<\varepsilon/2\)。从 \(D\) 选有限个中心，其 \(r/2\)-球覆盖 \(K\)，再选有理数 \(q_i\) 满足 \(|q_i-f(a_i)|<\varepsilon/2\)。上述分母为正，而每个非零权重都满足 \(d(x,a_i)<r\)。因此 \(g(x)\) 是若干与 \(f(x)\) 相差小于 \(\varepsilon\) 的数的凸组合，故 \(\|g-f\|_\infty<\varepsilon\)。
\end{proof}

这解释了为什么连续路径空间可以纳入第十二章的测度收敛理论。例如 \(C([0,1])\) 完备且可分，但它的闭单位球不紧；不同频率或不同窄峰的函数会破坏等度连续性。波兰性给出适宜的状态空间，紧族条件仍需另行证明。

\subsection{返回测度与局部分析}
\label{b1:subsec:A10503}

本附录的工具可以沿两条方向使用。若研究紧支集测试与局部积分，局部紧性提供紧邻域，闭集分离提供连续截断，可数基则进一步提供紧穷竭。若研究概率分布或随机路径，可分完备度量提供统一的可数逼近与序列极限，未必存在紧邻域，因此不能把 \(C_c\) 的方法原样搬过去。

在第\ref{b1:sec:1204}节证明单个概率测度的紧性时，可分性给出可数小球覆盖，再在每一尺度选有限子族；把各尺度的覆盖同时满足，得到全有界的闭集，完备性最后使它紧。现在读者可以清楚辨认这三个步骤：可数化、有限逼近、保留极限。它们不是同一种“紧性”的不同叫法，而是构造紧集时先后承担的不同任务。
